% ===== PRÉAMBULE CANONIQUE QCM — PHYSIQUE =====
% Sert à compiler controle.tex ET à la revalidation automatique du JSON
% (valider_qcm.py). NE PAS MODIFIER : packages figés.
% La bibliothèque TikZ "babel" est indispensable (conflit babel-french/circuitikz).
\documentclass[11pt,a4paper]{article}
\usepackage[utf8]{inputenc}\usepackage[T1]{fontenc}\usepackage{lmodern}
\usepackage[french]{babel}
\usepackage{amsmath,amssymb,mathtools}\usepackage{icomma}
\usepackage{siunitx}\sisetup{locale=FR}
\usepackage{xcolor}\usepackage{colortbl}
\usepackage{tikz}\usepackage{pgfplots}\pgfplotsset{compat=1.18}
\usepackage[siunitx,europeanresistors]{circuitikz}
\usepackage{enumitem}\usepackage{array,booktabs}
\usepackage[a4paper,margin=2.2cm]{geometry}
\usetikzlibrary{arrows.meta,calc,angles,quotes,positioning,patterns,%
  backgrounds,shapes.geometric,decorations.pathmorphing,%
  decorations.markings,intersections,babel}
\newcommand{\vv}[1]{\vec{#1}}\newcommand{\ds}{\displaystyle}
\newcommand{\R}{\mathbb{R}}\newcommand{\N}{\mathbb{N}}\newcommand{\Z}{\mathbb{Z}}
% ==============================================
\begin{document}
\noindent\textbf{PHYS-F04-Q01 --- Loi de gravitation : effet des masses et de la distance}\par
Deux corps $A$ et $B$ exercent l'un sur l'autre une force de gravitation d'intensité $F$. On multiplie la masse du corps $A$ par $8$ et la distance entre leurs centres par $2$. Par quel facteur l'intensité de la force de gravitation est-elle multipliée ?
\begin{enumerate}[label=\Alph*)]
\item $\times 4$
\item $\times 32$
\item $\times 2$
\item $\times \dfrac{1}{2}$
\end{enumerate}
\textit{Bonne réponse : C}\par
L'intensité obéit à $F=\mathcal{G}\,\dfrac{m_A m_B}{d^2}$. En multipliant $m_A$ par $8$ et $d$ par $2$, le nouveau facteur vaut $\dfrac{8}{2^2}=\dfrac{8}{4}=2$ : la force est multipliée par $2$.\par
\textit{Pièges :}
\begin{itemize}
\item A: loi supposée en $1/d$ au lieu de $1/d^2$ : $8/2=4$.
\item B: distance placée au numérateur (oubli du caractère inverse) : $8\times 2^2=32$.
\item D: rôles de la masse et de la distance intervertis : $2^2/8=1/2$.
\end{itemize}
\vspace{8pt}\hrule\vspace{8pt}
\noindent\textbf{PHYS-F04-Q02 --- Distance au centre contre altitude}\par
Au niveau du sol ($h=0$), la Terre exerce sur un satellite de masse $m$ une force d'attraction d'intensité $F_0$. On double la masse du satellite (elle devient $2m$) et on le place à l'altitude $h=R_T$, où $R_T$ est le rayon de la Terre. Quelle est la nouvelle intensité de la force d'attraction terrestre ?
\begin{enumerate}[label=\Alph*)]
\item $\dfrac{F_0}{2}$
\item $2F_0$
\item $\dfrac{F_0}{4}$
\item $F_0$
\end{enumerate}
\textit{Bonne réponse : A}\par
Au sol, $F_0=\mathcal{G}\,\dfrac{M_T m}{R_T^2}$. À l'altitude $h=R_T$, la distance au centre vaut $d=R_T+h=2R_T$ et la masse est $2m$ : $F=\mathcal{G}\,\dfrac{M_T (2m)}{(2R_T)^2}=\dfrac{2}{4}\,F_0=\dfrac{F_0}{2}$.\par
\textit{Pièges :}
\begin{itemize}
\item B: garde $d=R_T$ (confond altitude et distance au centre), seule la masse double : $2F_0$.
\item C: oublie le doublement de la masse : $1/2^2=F_0/4$.
\item D: oublie d'élever la distance au carré : $2/2=F_0$.
\end{itemize}
\vspace{8pt}\hrule\vspace{8pt}
\noindent\textbf{PHYS-F04-Q03 --- Variation de la pesanteur avec l'altitude}\par
On donne l'intensité de pesanteur au sol $g_0=\qty{9.8}{\metre\per\second\squared}$ et le rayon terrestre $R_T=\qty{6400}{\kilo\metre}$. Quelle est l'intensité de pesanteur $g$ à l'altitude $h=\qty{6400}{\kilo\metre}$ ?
\begin{enumerate}[label=\Alph*)]
\item \qty{4.90}{\metre\per\second\squared}
\item \qty{9.80}{\metre\per\second\squared}
\item \qty{1.09}{\metre\per\second\squared}
\item \qty{2.45}{\metre\per\second\squared}
\end{enumerate}
\textit{Bonne réponse : D}\par
$g=\dfrac{\mathcal{G}M_T}{(R_T+h)^2}=\dfrac{\mathcal{G}M_T}{(2R_T)^2}=\dfrac{g_0}{4}=\qty{2.45}{\metre\per\second\squared}$, car $R_T+h=2R_T$.\par
\textit{Pièges :}
\begin{itemize}
\item A: oubli du carré (loi en $1/d$) : $g_0/2=\qty{4.90}{\metre\per\second\squared}$.
\item B: confond altitude et distance au centre ($d=R_T$) : $g$ reste $g_0=\qty{9.80}{\metre\per\second\squared}$.
\item C: utilise le diamètre $2R_T$ au lieu du rayon ($d=3R_T$) : $g_0/9=\qty{1.09}{\metre\per\second\squared}$.
\end{itemize}
\vspace{8pt}\hrule\vspace{8pt}
\noindent\textbf{PHYS-F04-Q04 --- Pesanteur à la surface d'un autre astre}\par
Une planète a une masse égale à $2$ fois celle de la Terre et un rayon égal à $2$ fois celui de la Terre. Sachant que $g_0=\qty{9.8}{\metre\per\second\squared}$ à la surface de la Terre, quelle est l'intensité de pesanteur à la surface de cette planète ?
\begin{enumerate}[label=\Alph*)]
\item \qty{2.45}{\metre\per\second\squared}
\item \qty{4.90}{\metre\per\second\squared}
\item \qty{9.80}{\metre\per\second\squared}
\item \qty{19.6}{\metre\per\second\squared}
\end{enumerate}
\textit{Bonne réponse : B}\par
$g=\dfrac{\mathcal{G}(2M_T)}{(2R_T)^2}=\dfrac{2}{4}\,g_0=\dfrac{g_0}{2}=\qty{4.90}{\metre\per\second\squared}$ : le facteur $2$ de masse est plus que compensé par le carré du rayon doublé.\par
\textit{Pièges :}
\begin{itemize}
\item A: ne garde que l'effet du rayon (oubli du facteur de masse $2$) : $g_0/4=\qty{2.45}{\metre\per\second\squared}$.
\item C: rayon non élevé au carré : $2/2=g_0=\qty{9.80}{\metre\per\second\squared}$.
\item D: oubli de l'effet du rayon (garde seulement la masse $\times 2$) : $2g_0=\qty{19.6}{\metre\per\second\squared}$.
\end{itemize}
\vspace{8pt}\hrule\vspace{8pt}
\noindent\textbf{PHYS-F04-Q05 --- Poids apparent dans un ascenseur}\par
Une personne de masse $m=\qty{60}{\kilogram}$ est debout sur une balance, dans un ascenseur qui \textbf{monte en ralentissant}, avec une décélération de $\qty{2.0}{\metre\per\second\squared}$. On prend $g=\qty{10}{\metre\per\second\squared}$. Quelle valeur la balance indique-t-elle ?
\begin{enumerate}[label=\Alph*)]
\item \qty{720}{\newton}
\item \qty{600}{\newton}
\item \qty{480}{\newton}
\item \qty{120}{\newton}
\end{enumerate}
\textit{Bonne réponse : C}\par
L'ascenseur monte mais ralentit : le vecteur accélération est dirigé \emph{vers le bas}, donc $a=\qty{-2.0}{\metre\per\second\squared}$. La balance affiche le poids apparent $N=m(g+a)=60\times(10-2)=\qty{480}{\newton}$.\par
\textit{Pièges :}
\begin{itemize}
\item A: « monte » interprété comme accélération vers le haut : $N=m(g+a)=60\times 12=\qty{720}{\newton}$.
\item B: accélération ignorée, $N=mg=\qty{600}{\newton}$.
\item D: ne garde que $N=ma=60\times 2=\qty{120}{\newton}$, en oubliant le poids.
\end{itemize}
\vspace{8pt}\hrule\vspace{8pt}
\noindent\textbf{PHYS-F04-Q06 --- Poids apparent et impesanteur}\par
Un astronaute flotte à l'intérieur de la Station spatiale internationale, en orbite circulaire à environ $\qty{400}{\kilo\metre}$ d'altitude, où l'intensité de pesanteur vaut encore $g\approx\qty{8.7}{\metre\per\second\squared}$. Parmi les affirmations suivantes, laquelle est \textbf{correcte} ?
\begin{enumerate}[label=\Alph*)]
\item L'astronaute et la station sont en chute libre permanente : le poids apparent est nul, mais le poids réel (de l'ordre de $m\times\qty{8.7}{\metre\per\second\squared}$) ne l'est pas.
\item À cette altitude, la gravité terrestre est devenue nulle, ce qui provoque l'impesanteur.
\item Le poids réel de l'astronaute est nul, car aucun support ne s'exerce sur lui.
\item En orbite, la masse de l'astronaute a diminué, ce qui explique qu'il flotte.
\end{enumerate}
\textit{Bonne réponse : A}\par
En orbite, la gravité joue le rôle de force centripète : la station et l'astronaute sont en chute libre permanente. Le poids apparent (force de contact) est nul, mais le poids réel $mg$ ne l'est pas, puisque $g\approx\qty{8.7}{\metre\per\second\squared}$, soit près de \qty{90}{\percent} de sa valeur au sol. La masse, elle, est invariable.\par
\textit{Pièges :}
\begin{itemize}
\item B: idée fausse « plus de gravité en orbite » ; en réalité $g$ vaut encore $\approx\qty{8.7}{\metre\per\second\squared}$.
\item C: confusion poids réel / poids apparent : seul le poids apparent (contact) est nul.
\item D: confusion masse / poids : la masse est invariable, seul le poids apparent s'annule.
\end{itemize}
\vspace{8pt}\hrule\vspace{8pt}
\noindent\textbf{PHYS-F04-Q07 --- Vitesse d'un satellite en orbite circulaire}\par
Un satellite décrit une orbite circulaire à l'altitude $h=\qty{3600}{\kilo\metre}$. On donne le rayon terrestre $R_T=\qty{6400}{\kilo\metre}$ et $\mathcal{G}M_T=\num{4.0e14}$ (unités SI). Quelle est la vitesse orbitale du satellite ?
\begin{enumerate}[label=\Alph*)]
\item \qty{7.9e3}{\metre\per\second}
\item \qty{6.3e3}{\metre\per\second}
\item \qty{1.1e4}{\metre\per\second}
\item \qty{8.9e3}{\metre\per\second}
\end{enumerate}
\textit{Bonne réponse : B}\par
Le rayon de l'orbite est $r=R_T+h=\qty{1.0e7}{\metre}$. La vitesse vaut $v=\sqrt{\dfrac{\mathcal{G}M_T}{r}}=\sqrt{\dfrac{\num{4.0e14}}{\num{1.0e7}}}=\sqrt{\num{4.0e7}}\approx\qty{6.3e3}{\metre\per\second}$.\par
\textit{Pièges :}
\begin{itemize}
\item A: oubli de l'altitude ($r=R_T$) : $v=\sqrt{\mathcal{G}M_T/R_T}\approx\qty{7.9e3}{\metre\per\second}$.
\item C: utilise l'altitude seule comme rayon d'orbite ($r=h$) : $v\approx\qty{1.1e4}{\metre\per\second}$.
\item D: facteur $\sqrt{2}$ en trop (confusion avec la vitesse de libération) : $\approx\qty{8.9e3}{\metre\per\second}$.
\end{itemize}
\vspace{8pt}\hrule\vspace{8pt}
\noindent\textbf{PHYS-F04-Q08 --- Troisième loi de Kepler}\par
Deux satellites $A$ et $B$ décrivent des orbites circulaires autour de la Terre. Le rayon de l'orbite de $B$ est $4$ fois celui de $A$. D'après la 3\textsuperscript{e} loi de Kepler, que vaut le rapport $\dfrac{T_B}{T_A}$ des périodes de révolution ?
\begin{enumerate}[label=\Alph*)]
\item $4$
\item $64$
\item $2$
\item $8$
\end{enumerate}
\textit{Bonne réponse : D}\par
La 3\textsuperscript{e} loi de Kepler donne $T^2\propto r^3$, donc $\dfrac{T_B}{T_A}=\left(\dfrac{r_B}{r_A}\right)^{3/2}=4^{3/2}=(\sqrt{4})^3=2^3=8$.\par
\textit{Pièges :}
\begin{itemize}
\item A: relation supposée linéaire $T\propto r$ : $4^1=4$.
\item B: oubli de la racine carrée sur $T^2$ (relation lue $T\propto r^3$) : $4^3=64$.
\item C: relation supposée $T\propto\sqrt{r}$ : $4^{1/2}=2$.
\end{itemize}
\vspace{8pt}\hrule\vspace{8pt}
\noindent\textbf{PHYS-F04-Q09 --- Satellite géostationnaire}\par
Parmi les affirmations suivantes concernant un satellite \textbf{géostationnaire}, laquelle est \textbf{correcte} ?
\begin{enumerate}[label=\Alph*)]
\item Il a la même vitesse angulaire qu'un point situé sur l'équateur terrestre, mais une vitesse linéaire nettement plus grande.
\item Il a la même vitesse linéaire qu'un point de l'équateur, puisqu'ils effectuent un tour dans le même temps.
\item Sa période de révolution vaut environ $\qty{90}{\minute}$, comme celle de la Station spatiale internationale.
\item On peut le placer à n'importe quelle altitude, pourvu qu'il reste au-dessus de l'équateur.
\end{enumerate}
\textit{Bonne réponse : A}\par
Un satellite géostationnaire a la même période que la Terre ($T\approx\qty{24}{\hour}$) : il partage donc la vitesse angulaire $\omega$ d'un point de l'équateur. Mais comme $v=\omega r$ et que son rayon d'orbite est bien plus grand, sa vitesse linéaire est nettement supérieure. Cette période impose une altitude \emph{unique}, $\approx\qty{35800}{\kilo\metre}$.\par
\textit{Pièges :}
\begin{itemize}
\item B: confond vitesse angulaire et vitesse linéaire ; $v=\omega r$ dépend du rayon.
\item C: confond avec l'ISS ; le géostationnaire a $T\approx\qty{24}{\hour}$, pas $\qty{90}{\minute}$.
\item D: ignore l'altitude unique ($\approx\qty{35800}{\kilo\metre}$) imposée par $T\approx\qty{24}{\hour}$.
\end{itemize}
\vspace{8pt}\hrule\vspace{8pt}
\noindent\textbf{PHYS-F04-Q10 --- Énergie mécanique d'un satellite}\par
Un satellite en orbite circulaire possède une énergie cinétique $E_c=\qty{2.0e10}{\joule}$. Que vaut son énergie mécanique totale $E_m$ sur cette orbite ?
\begin{enumerate}[label=\Alph*)]
\item \qty{-4.0e10}{\joule}
\item \qty{0}{\joule}
\item \qty{-2.0e10}{\joule}
\item \qty{+2.0e10}{\joule}
\end{enumerate}
\textit{Bonne réponse : C}\par
Sur une orbite circulaire, $E_p=-2E_c$, donc $E_m=E_c+E_p=E_c-2E_c=-E_c=\qty{-2.0e10}{\joule}$. L'énergie mécanique d'un corps lié à la Terre est négative.\par
\textit{Pièges :}
\begin{itemize}
\item A: confond $E_m$ avec $E_p=-2E_c=\qty{-4.0e10}{\joule}$.
\item B: utilise $E_p=-E_c$ au lieu de $-2E_c$, d'où $E_m=0$.
\item D: oubli du signe négatif : $E_m=+E_c=\qty{+2.0e10}{\joule}$.
\end{itemize}
\vspace{8pt}\hrule\vspace{8pt}
\end{document}
